\section*{Original definitions and computational models (pp.9:1--2,9:11--14)}
A polytope is the convex hull of a finite set of points in Euclidean space. Equivalently, it is a compact set that is defined by finitely many linear inequalities.
A shadow of a polytope $P\subseteq\mathbb R^n$ is a set of the form $L(P)$, where $L:\mathbb R^n\to\mathbb R^2$ is a linear map. The shadows of $P$ are two-dimensional polygons, and hence typically much simpler than $P$. The shadow complexity of $P$ is
\[
\sigma(P)=\max_L |\vertset(L(P))|,
\]
where $L$ is a linear map and $\vertset(Q)$ is the vertex set of the polytope $Q$.
Let $\mathbb F$ be a field. For a list of variables $x=(x_1,\ldots,x_n)$ and $\alpha\in\mathbb N^n$, let $x^\alpha$ be the monomial $\prod_{i=1}^n x_i^{\alpha_i}$. A polynomial $f\in\mathbb F[x_1,\ldots,x_n]$ is a formal sum of the form $\sum_{\alpha\in\mathbb N^n}a_\alpha x^\alpha$ where $\supp(f):=\{\alpha\in\mathbb N^n:a_\alpha\ne0\}$ is finite. The Newton polytope of $f$ is
\[\Newt(f):=\conv(\supp(f)),\]
where $\conv(\cdot)$ denotes the convex hull.
A polynomial $f$ is transparent if $|\supp(f)|=\sigma(\Newt(f))$. In other words, there is a linear map $L:\mathbb R^n\to\mathbb R^2$ which maps $\supp(f)$ to distinct convexly independent points in $\mathbb R^2$.

As the standard model of computation of polynomials we take the arithmetic circuit model. It is a (finite) directed acyclic graph whose every node has in-degree zero or two. Nodes of in-degree zero (input nodes) are labelled with a constant or a variable. Nodes of in-degree two are labelled with operations $+$ or $\times$. Every node of the circuit computes a polynomial in $\mathbb F$ in the natural way. As the size of the circuit, we take the number of nodes. A circuit is called a formula if its underlying graph is a tree. For more background and motivation, see [44] and references within.
Our focus is mainly on monotone computation. A polynomial over $\mathbb R$ is monotone if all of its coefficients are non-negative. Similarly, a monotone arithmetic circuit can use only non-negative constants.

In a Laurent polynomial, variables are allowed to have negative integer exponents. The notions of $\supp(f)$ and Newton polytope of $f$ are defined in the obvious manner. A Laurent polynomial over $\mathbb R$ is monotone, if all of its coefficients are non-negative.
A high powered circuit (h.p.-circuit for short) is an arithmetic circuit in which every input node is labelled by a term $ax_1^{k_1}\cdots x_n^{k_n}$ with $a\in\mathbb F$ and $k_1,\ldots,k_n\in\mathbb Z$. The size of the circuit is the number of its nodes.
In other words, we have given the circuit a power to compute every term $ax^\alpha$ at a unit cost. This is especially important in the case of h.p.-formula. An arithmetic formula of size $s$ can compute a polynomial of degree at most $s$, whereas there is no such restriction in an h.p.-formula. Furthermore, we have allowed the variables to have negative exponents and hence an h.p.-circuit computes a Laurent polynomial instead of a polynomial. But this is only a cosmetic detail.

A high power projection (h.p.-projection for short) is a homomorphism
\[\pi:\mathbb F[x_1,\ldots,x_n]\to\mathbb F[y_1,\ldots,y_m,y_1^{-1},\ldots,y_m^{-1}]\]
such that $\pi(x_i)=a_i y^{\alpha_i}$ for every $x_i$, where $a_i\in\mathbb F$ and $\alpha_i\in\mathbb Z^m$, and for every $f\in\mathbb F[x_1,\ldots,x_n]$,
\[\pi(f(x_1,\ldots,x_n))=f(\pi(x_1),\ldots,\pi(x_n)).\]
The constants $a_i$ are called the coefficients of $\pi$ and $\alpha_i$ its exponents. If $\mathbb F=\mathbb R$ and $\pi$ has non-negative coefficients, then $\pi$ is called monotone.
An h.p.-projection $\pi$ induces a linear map $L_\pi:\mathbb R^n\to\mathbb R^m$ by setting $L_\pi(e_i)=\alpha_i$ and extending it linearly to $\mathbb R^n$. It follows that $\supp(\pi(f))\subseteq L_\pi(\supp(f))$. The inclusion may be strict, as some monomials of $f$ can cancel out in the projection. If no cancellations occur, we indeed have $\Newt(\pi(f))=L_\pi(\Newt(f))$. This is satisfied, e.g., if $f$ is monotone and the coefficients of $\pi$ are positive, or if the coefficients are algebraically independent.
